% page 608 du fac-simile (image p-607 du scan)
% bloc 154.9 x 246.3 mm ; marge gauche 8.0 mm
\pgc{-12.700mm}{0.000mm}{
\l{\cel{2}600}
\l{}
\l{tumultar.\cel{2}(\sou{netrans}.) (\sou{aludante personi asemblita}). Agitar}
\l{\cel{3}su, bruisoze, sen-ordine, generale pro granda nekontent-}
\l{\cel{3}eso e kom ensemblo. - (\sou{anke metaf.}) - DEFIS.}
\l{}
\l{\sou{tuno}. Mezur-unajo (1.000 kilogrami) por evaluar la kargajo}
\l{\cel{3}tota di la vari sur navo, sur vagono, sur kamiono auto-}
\l{\cel{3}mobila. - DEFIRS.}
\l{}
\l{\sou{tunelo}. (\sou{tekn.}) Paseyo subsula, facita sub rivero; sub flu-}
\l{\cel{3}vio; sub parto de maro, qua quaze,klemesas inter du tero-}
\l{\cel{3}langi ; sub monto ; tra monto, e c. - DEFIRS.}
\l{}
\l{\sou{tuniko.}\cel{2}I. (\sou{epoki antiqua}).\cel{2}Subvesto. - II.(\sou{armeo}). Redin-}
\l{\cel{3}goto kurta di uniformo. - III. Nomo di ula peli tenua,di}
\l{\cel{3}ula gaini membranatra (ex.: la tuniki di la okulo, di la}
\l{\cel{3}arterii, di la onyono). - DEFIRS.}
\l{}
\l{\sou{tupio}. Ludilo puerala, konatra, quan onu rotacigas sur lua}
\l{\cel{3}pinto fera ; onu startas lu per kordeto spulita an-cirkume}
\l{\cel{3}- e quan onu desspulas tre rapide, per bruska tiro di}
\l{\cel{3}la kordeto.- F.}
\l{}
\l{\sou{turo}. Voyajo segun parkuro plu o min cirklokurvatra, maxim-}
\l{\cel{3}multa-kaze por sua plezuro. - DEFR.}
\l{}
\l{\sou{turbo}. Quanto tre granda de personi ordinara. FIS.}
\l{}
\l{\sou{turbano}. I. Kapo-vesto virala, che la Orientani - tuko spu-}
\l{\cel{3}lita cirkum la kapo. - II. Kapo-vesto mulierala por la}
\l{\cel{3}bali, qua imitas la turbano virala. - DEFIRS.}
\l{}
\l{\sou{turbino}. (\sou{tekn.)}\cel{3}Roto trogoza od helicoza, qua kaptas la}
\l{\cel{3}energio de aquo fluanta o falanta, de vaporo spricanta,}
\l{\cel{3}de gaso, e qua transmisas ta energio per ax-arboro. -}
\l{\cel{3}DEFIRS.}
\l{}
\l{\sou{turboto}. (\sou{zool.}) Marfisho, ek la familio \textquotedbl{}pleuronekti\textquotedbl{}, di}
\l{\cel{3}qua la karno saporas delikate. - L.\cel{1}\sou{rhombus maximus}.}
\l{\cel{3}- DEFR.}
\l{}
\l{\sou{turdo}. (\sou{zool.})\cel{2}Ucelo ek la klaso \textquotedbl{}paseri dentirostra\textquotedbl{}, qua}
\l{\cel{3}formacas speco ek la genero \textquotedbl{}merlo\textquotedbl{}, kun plumaro blanka}
\l{\cel{3}e bruna. - L.\cel{1}\sou{turdus}. -}
\l{}
\l{\sou{turgecar}. (\sou{netrans.)}\cel{2}(\sou{fiziol.}) (\sou{aludante ula organi, preci-}}
\l{\cel{3}\sou{pue le sexuala})\cel{2}Inflesar e rigideskar da la adfluado di}
\l{\cel{3}sango od altra fluidi. - DEFIS.}
\l{}
\l{\sou{turkezo}. Lapido matida, bele blua. - DEFIS.}
\l{}
\l{\sou{turmo}. I. Konstrukturo cilindra o plur-edra, tre alta kompare}
\l{\cel{3}a la dimensioni di lua bazo, sola, o salianta ek edifici}
\l{\cel{3}vicina, e qua dominacas urbo, kastelo. (ex.: la turmo da}
\l{\cel{3}Eiffel, en Paris). - II. (\sou{shakoludo}).Piono turmo-forma,}
\l{\cel{3}en la angulo di la shako-plako. - D.}
}
